\chapter{التحريكيات في بعدين}\label{ch:g12:kinematics-2d}

راقب سيارة تدور في دوّار \omterm{def:g12:kinematics-2d:velocity}{بسرعة} ثابتة \qty{30}{km/h}: فلا تتحرك
الإبرة قط، ومع ذلك يحسّ كل راكب بجذب جانبي --- فالحركة
تتغير، وليست \omterm{def:g12:kinematics-2d:velocity}{السرعة} هي المتغيرة. ولم يعد عدد واحد لكل لحظة
يكفي: فالحركة في المستوي تحتاج سهامًا. ويزوّد هذا الفصل
نقطةً متحركة بثلاثة --- الموضع \omterm{def:g12:kinematics-2d:velocity}{والسرعة} والتسارع ---
مسلسَلةً بالاشتقاق، ويقرأ عائلات كاملة من الحركات من
هندستها.

\section{متجهة الموضع}

لا وجود للحركة إلا بالنسبة إلى \omterm{def:g10:relative-motion:frame}{مرجع} مذكور --- أي جسم صلب
مع ساعة (\cref{ch:g10:relative-motion}). والجديد هذه السنة هو
الدقة: أرفِق \omterm{def:g10:relative-motion:frame}{بالمرجع} مبدأً $O$ ومحورين متعامدين
مدرّجين $Ox$ و$Oy$، وسجّل أين تقع النقطة في كل
لحظة، زوجَ دالّتين.

\begin{definition}[متجهة الموضع]\label{def:g12:kinematics-2d:position}
في \omterm{def:g10:relative-motion:frame}{مرجع} مزوَّد بمبدأ $O$ ومحورين $Ox$ و$Oy$، تكون
\emph{متجهة الموضع}\index{متجهة الموضع} لنقطة متحركة $M$
عند اللحظة $t$ هي $\vect{OM}(t)$، وإحداثياها
$\bigl(x(t),\, y(t)\bigr)$: أي دالّتان للزمن، واحدة لكل محور.
والمنحني الذي يرسمه $M$ بتغير $t$ هو مسارها.
\end{definition}

\section{متجهة السرعة}

\begin{definition}[متجهة السرعة]\label{def:g12:kinematics-2d:velocity}
على مدى فترة قصيرة تنتقل النقطة بمقدار
$\vect{OM}(t+\Delta t) - \vect{OM}(t)$؛ فاقسم على $\Delta t$ ودع
الفترة تتقلص: وهذا هو الاشتقاق من درس الرياضيات، مطبَّقًا
على كل إحداثي. و\emph{متجهة سرعة}\index{متجهة السرعة}
النقطة $M$ هي مشتقة متجهة موضعها،
$\vect v(t) = \bigl(x'(t),\, y'(t)\bigr)$، \omterm{def:g10:orders-of-magnitude:unit}{بوحدة}
\unit{m/s}. وناظمها $v = \sqrt{x'^2 + y'^2}$ هو
\emph{السرعة}\index{سرعة} --- أي العدد الوحيد الذي يعرضه عدّاد السرعة.
\end{definition}

\begin{example}[طائرة مسيّرة في طيران أفقي]\label{ex:g12:kinematics-2d:drone}
تطير مسيّرة بحيث $x(t) = 4.0\,t$ و$y(t) = 3.0\,t$ (\omterm{def:g10:orders-of-magnitude:unit}{بالمتر} و
\omterm{def:g10:orders-of-magnitude:unit}{الثانية}): فمسارها المستقيم $y = \tfrac34 x$، وسرعتها
$\vect v = (4.0, 3.0)$ في كل حين، وقيمة \omterm{def:g12:kinematics-2d:velocity}{السرعة}
$\sqrt{4.0^2 + 3.0^2} = \qty{5.0}{m/s}$، وهي ثابتة.
\end{example}

\begin{proposition}[مماسّة للمسار]\label{prop:g12:kinematics-2d:tangent}
في كل لحظة، تكون $\vect v(t)$ مماسّة \omterm{def:g10:relative-motion:trajectory}{للمسار} عند
الموضع الحالي للنقطة $M$، وتتجه في جهة السير.
\end{proposition}

\begin{proof}
الانتقال وتر من أوتار \omterm{def:g10:relative-motion:trajectory}{المسار}؛ وكلما تقلص $\Delta t$،
صار اتجاه الوتر اتجاه المماسّ.
\end{proof}

\section{متجهة التسارع}

\begin{definition}[متجهة التسارع]\label{def:g12:kinematics-2d:acceleration}
\emph{متجهة تسارع}\index{متجهة التسارع} النقطة $M$ هي
مشتقة متجهة سرعتها،
$\vect a(t) = \bigl(v_x'(t),\, v_y'(t)\bigr) =
\bigl(x''(t),\, y''(t)\bigr)$، \omterm{def:g10:orders-of-magnitude:unit}{بوحدة} \unit{m/s^2}: أي كم، وفي
أي اتجاه، تتغير \omterm{def:g12:kinematics-2d:velocity}{متجهة السرعة} في اللحظة الحالية.
\end{definition}

\begin{remark}[تسارع بلا ازدياد في السرعة]\label{rem:g12:kinematics-2d:turning}
يكون $\vect a \neq \vect 0$ كلما تغيرت \emph{متجهة} \omterm{def:g12:kinematics-2d:velocity}{السرعة}
--- في الطول \emph{أو في الاتجاه}: فالسيارة الكابحة تتسارع
(إلى الوراء)؛ والسيارة المنعطفة \omterm{def:g12:kinematics-2d:velocity}{بسرعة} ثابتة تتسارع (جانبيًّا).
وهذا هو المقدار الذي تتحكم فيه القوى
(\cref{ch:g11:forces-and-motion,ch:g12:newtons-laws}).
\end{remark}

\begin{method}[قراءة تصوير متعاقب]\label{met:g12:kinematics-2d:chrono}
\emph{التصوير المتعاقب}\index{تصوير متعاقب} يسجّل
مواضع $M_0, M_1, M_2, \dots$ نقطةٍ متحركة على فترات زمنية
متساوية $\Delta t$.
\begin{enumerate}
  \item تباعد النقاط يقرأ \omterm{def:g12:kinematics-2d:velocity}{السرعة}: فالتباعدات المتساوية تعني \omterm{def:g12:kinematics-2d:velocity}{سرعة} ثابتة؛
        والمتسعة تعني تسارعًا؛ والمتضايقة تعني تباطؤًا.
  \item \omterm{def:g12:kinematics-2d:velocity}{والسرعة} عند $M_i$ هي الوتر المتناظر،
        $\vect v_i \approx \vect{M_{i-1}M_{i+1}}/(2\,\Delta t)$،
        مرسومًا عند $M_i$، مماسًّا \omterm{def:g10:relative-motion:trajectory}{للمسار} المنقَّط.
  \item وللتسارع عند $M_i$: انسخ $\vect v_{i-1}$ و
        $\vect v_{i+1}$ إلى $M_i$، وارسم
        $\Delta\vect v = \vect v_{i+1} - \vect v_{i-1}$ رأسًا إلى رأس؛
        ومنه $\vect a_i \approx \Delta\vect v/(2\,\Delta t)$.
\end{enumerate}
\end{method}

\begin{omfigure}
\begin{tikzpicture}[scale=1.05, font=\small]
  \fill (0,2.05) circle (2pt) (0.8,1.95) circle (2pt)
    (1.6,1.72) circle (2pt) (2.4,1.36) circle (2pt)
    (3.2,0.87) circle (2pt) (4.0,0.25) circle (2pt);
  \node[above] at (0,2.05) {$M_0$};
  \node[above] at (1.6,1.72) {$M_2$};
  \node[above right] at (2.4,1.36) {$M_3$};
  \node[above right] at (3.2,0.87) {$M_4$};
  \draw[-Stealth, omDef, thick] (1.6,1.72) -- ++(1.2,-0.44)
    node[right] {$\vect v_2$};
  \draw[-Stealth, omDef, thick] (3.2,0.87) -- ++(1.2,-0.83)
    node[right] {$\vect v_4$};
  \draw[-Stealth, omDef, dashed] (2.4,1.36) -- ++(1.2,-0.44);
  \draw[-Stealth, omDef] (2.4,1.36) -- ++(1.2,-0.83);
  \draw[-Stealth, omThm, very thick] (3.6,0.92) -- ++(0,-0.39)
    node[right] {$\Delta\vect v$};
  \draw[-Stealth, omThm, very thick] (2.4,1.36) -- ++(0,-0.75)
    node[left] {$\vect a$};
\end{tikzpicture}

{\small \omterm{met:g12:kinematics-2d:chrono}{تصوير متعاقب} لكرة مقذوفة: تُنسخ $\vect v_2$ و$\vect
v_4$، المبنيتان من الأوتار، إلى $M_3$؛ ويسلّم فرقهما
$\Delta\vect v$ المتجهةَ $\vect a$ --- عمودية إلى أسفل.}
\end{omfigure}

\section{الحركات المستقيمة}

على مستقيم يكفي محور واحد: $x(t)$ و$v(t) = x'(t)$ و
$a(t) = v'(t)$، وكلها أعداد ذات \omterm{def:g10:signals-and-waves:signal}{إشارة}.

\begin{definition}[الحركة المستقيمة المنتظمة]\label{def:g12:kinematics-2d:uniform}
تكون الحركة \emph{مستقيمة منتظمة}\index{حركة مستقيمة
منتظمة} إذا كانت $\vect v$ متجهةً ثابتة: \omterm{def:g10:relative-motion:trajectory}{فالمسار} مستقيم،
\omterm{def:g12:kinematics-2d:velocity}{والسرعة} ثابتة، و$\vect a = \vect 0$. ومنه $x(t) = v\,t + x_0$.
\end{definition}

\begin{proposition}[الحركة المستقيمة المتسارعة بانتظام]\label{prop:g12:kinematics-2d:uarm}
\index{حركة متسارعة بانتظام}%
إذا تحركت نقطة على $Ox$ بتسارع \emph{ثابت} $a$،
منطلقةً عند $t = 0$ من الموضع $x_0$ \omterm{def:g12:kinematics-2d:velocity}{بسرعة} $v_0$، فإن
\[
v(t) = a\,t + v_0,
\qquad
x(t) = \tfrac12\,a\,t^2 + v_0\,t + x_0 .
\]
\end{proposition}

\begin{proof}
$v$ دالة أصلية للثابت $a$: أي $v(t) = a t + C$، ويفرض
$t = 0$ أن $C = v_0$. وبدوره $x$ دالة أصلية \omterm{def:g12:kinematics-2d:velocity}{للسرعة}
$a t + v_0$: أي $x(t) = \tfrac12 a t^2 + v_0 t + C'$، مع $C' = x_0$.
\end{proof}

\begin{omfigure}
\begin{tikzpicture}[font=\small]
\begin{axis}[omaxis, name=ax1, width=8.4cm, height=3.2cm,
    ylabel={$x$ (\unit{m})}, xmin=0, xmax=4, ymin=0, ymax=21]
  \addplot[omDef, thick, domain=0:4, samples=40] {0.5*2*x^2 + x};
\end{axis}
\begin{axis}[omaxis, name=ax2, at={(ax1.below south west)},
    anchor=north west, yshift=-3mm, width=8.4cm, height=3.2cm,
    ylabel={$v$ (\unit{m/s})}, xmin=0, xmax=4, ymin=0, ymax=10]
  \addplot[omThm, thick, domain=0:4, samples=2] {2*x + 1};
\end{axis}
\begin{axis}[omaxis, at={(ax2.below south west)}, anchor=north west,
    yshift=-3mm, width=8.4cm, height=3.2cm, xlabel={$t$ (\unit{s})},
    ylabel={$a$ (\unit{m/s^2})}, xmin=0, xmax=4, ymin=0, ymax=3]
  \addplot[omProp, thick, domain=0:4, samples=2] {2};
\end{axis}
\end{tikzpicture}

{\small \omterm{prop:g12:kinematics-2d:uarm}{حركة متسارعة بانتظام} ($a = \qty{2.0}{m/s^2}$،
$v_0 = \qty{1.0}{m/s}$، $x_0 = 0$): كل مخطط مشتقة
الذي فوقه --- قطع مكافئ، ثم ميله، ثم ميل ميله.}
\end{omfigure}

\begin{example}[مسافة التوقف]\label{ex:g12:kinematics-2d:braking}
سيارة \omterm{def:g12:kinematics-2d:velocity}{بسرعة} $v_0$ تكبح بتباطؤ ثابت $a$
(أي بتسارع $-a$). فتتوقف حين $v(t) = -a t + v_0 = 0$، عند
$t_s = v_0/a$، وقد قطعت $d = -\tfrac12 a t_s^2 + v_0 t_s =
v_0^2/(2a)$. والتربيع هو عنوان السلامة على الطرق:
فمضاعفة \omterm{def:g12:kinematics-2d:velocity}{السرعة} \emph{تربّع} مسافة التوقف --- فبأخذ $a =
\qty{6.0}{m/s^2}$، تكون \qty{16}{m} من \qty{50}{km/h}، و\qty{64}{m} من
\qty{100}{km/h}.
\end{example}

\section{الحركة الدائرية المنتظمة ومعلم فريني}

\begin{definition}[الحركة الدائرية المنتظمة]\label{def:g12:kinematics-2d:ucm}
تكون الحركة \emph{دائرية منتظمة}\index{حركة دائرية منتظمة}
إذا كان \omterm{def:g10:relative-motion:trajectory}{المسار} دائرة نصف قطرها $R$ \omterm{def:g12:kinematics-2d:velocity}{والسرعة} $v$
ثابتة. و\emph{الدور}\index{دور!الدوران} $T$ هو
مدة دورة واحدة، و\emph{التواتر}\index{تواتر!الدوران}
$f = 1/T$ (\omterm{def:g10:orders-of-magnitude:unit}{بوحدة} \unit{Hz}) هو عدد الدورات في \omterm{def:g10:orders-of-magnitude:unit}{الثانية}:
$v = 2\pi R/T = 2\pi R f$.
\end{definition}

\omterm{def:g12:kinematics-2d:velocity}{سرعة} ثابتة --- ومع ذلك تدور \omterm{def:g12:kinematics-2d:velocity}{متجهة السرعة}: أي أن ثمة \emph{تسارعًا}
فعلًا. فإلى أي جهة، وبأي مقدار؟

\begin{theorem}[التسارع المركزي]\label{thm:g12:kinematics-2d:centripetal}
\index{تسارع مركزي}%
في \omterm{def:g12:kinematics-2d:ucm}{حركة دائرية منتظمة} نصف قطرها $R$ وسرعتها $v$، يتجه
التسارع --- ويسمّى \emph{مركزيًّا} --- من النقطة
المتحركة نحو مركز الدائرة، وناظمه ثابت
\[
a = \frac{v^2}{R}.
\]
\end{theorem}

\begin{proof}
ضع مبدأ المحاور في مركز الدائرة واجعل $\omega = v/R$، بحيث يقابل
القوس $v t$ الزاوية $\omega t$: فيكون الموضع
$x(t) = R\cos(\omega t)$، $y(t) = R\sin(\omega t)$. وبالاشتقاق،
$\vect v = (-R\omega\sin(\omega t),\, R\omega\cos(\omega t))$، وناظمها
ثابت $R\omega = v$؛ وبالاشتقاق مرة أخرى،
$\vect a = -\omega^2\,\vect{OM}$: أي معاكسة \omterm{def:g12:kinematics-2d:position}{لمتجهة الموضع} ---
نحو المركز --- وناظمها $\omega^2 R = v^2/R$.
\end{proof}

\begin{omfigure}
\begin{tikzpicture}[scale=1.0, font=\small]
  \draw[thick] (0,0) circle (1.7);
  \fill (0,0) circle (1.5pt); \node[below left] at (0,0) {$O$};
  \draw[dashed, omExo] (0,0) -- (215:1.7)
    node[midway, below right] {$R$};
  \fill (35:1.7) circle (2pt); \node[above right] at (35:1.7) {$M$};
  \draw[-Stealth, omDef, thick] (35:1.7) -- ++(-0.69,0.98)
    node[above] {$\vect v$};
  \draw[-Stealth, omThm, very thick] (35:1.7) -- ++(215:0.95)
    node[below right] {$\vect a$};
\end{tikzpicture}

{\small \omterm{def:g12:kinematics-2d:ucm}{حركة دائرية منتظمة}: $\vect v$ مماسّة وطولها ثابت؛
و$\vect a$ مصوَّبة نحو المركز وطولها ثابت $v^2/R$ --- فتدير
\omterm{def:g12:kinematics-2d:velocity}{السرعة} أبدًا دون أن تمدّها قط.}
\end{omfigure}

\begin{definition}[معلم فريني]\label{def:g12:kinematics-2d:frenet}
\emph{معلم فريني}\index{معلم فريني} عند الموضع الحالي
لنقطة على \omterm{def:g10:relative-motion:trajectory}{مسار} منحنٍ هو زوج المتجهتين الواحديتين
المتعامدتين $(\vect{u_t}, \vect{u_n})$: فالمتجهة $\vect{u_t}$ مماسّة، في جهة
الحركة؛ و$\vect{u_n}$ ناظمية، نحو مركز الدائرة ذات
نصف القطر $R$ التي تلائم المنحني هناك أفضل ملاءمة.
\end{definition}

\begin{proposition}[التسارع في معلم فريني]\label{prop:g12:kinematics-2d:frenetacc}
على أي \omterm{def:g10:relative-motion:trajectory}{مسار} مستوٍ نصف قطره المحلي $R$، يُقطع \omterm{def:g12:kinematics-2d:velocity}{بسرعة} $v(t)$، يكون
\[
\vect a = a_t\,\vect{u_t} + a_n\,\vect{u_n},
\qquad a_t = \frac{\dd v}{\dd t},
\qquad a_n = \frac{v^2}{R}:
\]
و\emph{التسارع المماسّي}\index{تسارع مماسّي}
$a_t$ يغيّر \omterm{def:g12:kinematics-2d:velocity}{السرعة}؛ و\emph{التسارع
الناظمي}\index{تسارع ناظمي} $a_n$ يدير \omterm{def:g12:kinematics-2d:velocity}{متجهة السرعة}.
\end{proposition}
\admitted

\begin{remark}[أين يقيم البرهان]\label{rem:g12:kinematics-2d:frenetproof}
لقد برهنّا على الطرفين: $a_t$ خالصًا على مستقيم
(\cref{prop:g12:kinematics-2d:uarm})، و$a_n$ خالصًا على دائرة
(\cref{thm:g12:kinematics-2d:centripetal})؛ أما التفكيك
العام فيُشتقّ في مجلد السنة الأولى.
\end{remark}

\begin{example}[الكبح داخل منعطف]\label{ex:g12:kinematics-2d:bend}
تدخل سيارة منعطفًا نصف قطره $R = \qty{50}{m}$ \omterm{def:g12:kinematics-2d:velocity}{بسرعة} $v = \qty{15}{m/s}$
وهي تكبح بمقدار $\dd v/\dd t = \qty{-3.0}{m/s^2}$: فيكون $a_t =
\qty{-3.0}{m/s^2}$، $a_n = 15^2/50 = \qty{4.5}{m/s^2}$، ومنه $a =
\sqrt{3.0^2 + 4.5^2} \approx \qty{5.4}{m/s^2}$، متجهًا إلى الوراء
\emph{وإلى} داخل المنعطف. ويجب أن توفّر الإطارات الاثنين معًا ---
ولذلك يكبح سائقو السباقات \emph{قبل} المنعطف.
\end{example}

\section{تمارين}

\begin{exercise}[$\star$]\label{exo:g12:kinematics-2d:1}
من أجل $x(t) = 3.0\,t$، $y(t) = 4.0\,t$ (\omterm{def:g10:orders-of-magnitude:unit}{بالمتر} \omterm{def:g10:orders-of-magnitude:unit}{والثانية})، أعطِ
$\vect v$، وقيمة \omterm{def:g12:kinematics-2d:velocity}{السرعة}، \omterm{def:g10:relative-motion:trajectory}{والمسار}، وصنف الحركة.
\end{exercise}

\begin{exercise}[$\star$]\label{exo:g12:kinematics-2d:2}
على \omterm{met:g12:kinematics-2d:chrono}{تصوير متعاقب} يُلتقط كل \qty{0.10}{s}، تقع النقاط على
مستقيم، ومتباعدة بانتظام بمقدار \qty{0.35}{m}. عيّن الحركة و
احسب \omterm{def:g12:kinematics-2d:velocity}{السرعة}. وعلامَ كانت تدلّ تباعدات متسعة؟
\end{exercise}

\begin{exercise}[$\star$]\label{exo:g12:kinematics-2d:3}
من أجل $x(t) = 2.0\,t^2$ على $Ox$ (\omterm{def:g10:orders-of-magnitude:unit}{بالمتر} \omterm{def:g10:orders-of-magnitude:unit}{والثانية})، اشتقّ $v(t)$
و$a(t)$، وسمِّ صنف الحركة، واحسب \omterm{def:g12:kinematics-2d:velocity}{السرعة} عند
$t = \qty{3.0}{s}$.
\end{exercise}

\begin{exercise}[$\star$]\label{exo:g12:kinematics-2d:4}
يركب \omterm{def:g11:work-of-force:power}{حصان} دوّامة دائرةً نصف قطرها \qty{4.5}{m}، ودورة كل
\qty{8.0}{s}: احسب \omterm{def:g12:kinematics-2d:ucm}{التواتر} \omterm{def:g12:kinematics-2d:velocity}{والسرعة} والتسارع
(ناظمًا واتجاهًا).
\end{exercise}

\begin{exercise}[$\star$]\label{exo:g12:kinematics-2d:5}
يصعد مخطط $v(t)$ مستقيمًا من \qty{2.0}{m/s} عند $t = 0$ إلى
\qty{10.0}{m/s} عند $t = \qty{4.0}{s}$. اقرأ التسارع،
ثم المسافة المقطوعة (وهي المساحة تحت المخطط).
\end{exercise}

\begin{exercise}[$\star\star$]\label{exo:g12:kinematics-2d:6}
يأخذ قطار سريع منعطفًا نصف قطره \qty{6.0}{km} \omterm{def:g12:kinematics-2d:velocity}{بسرعة}
\qty{320}{km/h}. احسب $a_n$ وقارنه بالمقدار $g$؛ ولماذا يجب أن
تتجنب الخطوط السريعة المنعطفات الضيقة، ولو كانت مائلة ميلًا تامًّا؟
\end{exercise}

\begin{exercise}[$\star\star$]\label{exo:g12:kinematics-2d:7}
لعربة على سكة مستقيمة $x(t) = -1.2\,t^2 + 8.0\,t + 3.0$
(\omterm{def:g10:orders-of-magnitude:unit}{بالمتر} \omterm{def:g10:orders-of-magnitude:unit}{والثانية}). جِد $v(t)$ و$a$، ولحظة الانعكاس وموضعه،
وصِف الحركة قبله وبعده.
\end{exercise}

\begin{exercise}[$\star\star$]\label{exo:g12:kinematics-2d:8}
\omterm{met:g12:kinematics-2d:chrono}{تصوير متعاقب} على مستقيم، كل $\Delta t = \qty{0.20}{s}$،
يعطي $x = 0$، $0.10$، $0.28$، $0.54$، $\qty{0.88}{m}$. قدّر
\omterm{def:g12:kinematics-2d:velocity}{السرعة} عند النقاط الداخلية الثلاث
(\cref{met:g12:kinematics-2d:chrono})، ثم التسارع؛ واستنتج.
\end{exercise}

\begin{exercise}[$\star\star$]\label{exo:g12:kinematics-2d:9}
يدور القمر حول الأرض على شبه دائرة نصف قطرها \qty{3.84e8}{m}
في \qty{27.3}{d}. احسب سرعته وتسارعه المركزي؛
وقارن الأخير بالمقدار $g$، علمًا أن القمر يبعد نحو $60$ نصف قطر
أرضي (\cref{ch:g10:universal-gravitation}).
\end{exercise}

\begin{exercise}[$\star\star$]\label{exo:g12:kinematics-2d:10}
حلّة غسّالة نصف قطرها \qty{25}{cm} تدور \omterm{def:g12:kinematics-2d:velocity}{بسرعة} $1200$
دورة في الدقيقة. احسب $f$ و$T$، \omterm{def:g12:kinematics-2d:velocity}{وسرعة} الحافة وتسارع
الحافة بمضاعفات $g$. ولماذا يبقى الغسيل ملتصقًا
بينما يغادر الماء؟
\end{exercise}

\begin{exercise}[$\star\star$]\label{exo:g12:kinematics-2d:11}
سيارة \omterm{def:g12:kinematics-2d:velocity}{بسرعة} \qty{90}{km/h} تكبح بمقدار ثابت $a = \qty{7.5}{m/s^2}$:
احسب زمن التوقف ومسافة التوقف، وارسم $v(t)$.
\end{exercise}

\begin{exercise}[$\star\star\star$]\label{exo:g12:kinematics-2d:12}
من أجل $x(t) = R\cos(\omega t)$، $y(t) = R\sin(\omega t)$: احسب
$\vect v$؛ وتحقّق من أن \omterm{def:g12:kinematics-2d:velocity}{السرعة} هي الثابت $R\omega$ ومن أن
$\vect v \perp \vect{OM}$؛ واحسب $\vect a$، وبيّن أنها
$-\omega^2\,\vect{OM}$، واستخرج $a = v^2/R$.
\end{exercise}

\begin{exercise}[$\star\star\star$]\label{exo:g12:kinematics-2d:13}
تدخل دراجة نارية منعطفًا نصف قطره \qty{80}{m} \omterm{def:g12:kinematics-2d:velocity}{بسرعة} \qty{20}{m/s}،
وهي تكبح بمقدار \qty{2.5}{m/s^2}. احسب $a_t$ و$a_n$، \omterm{def:g10:refraction:angles}{وناظم}
$\vect a$ وزاويتها مع جهة الحركة. وتتماسك الإطارات
حتى \qty{7.0}{m/s^2} في المجموع: فهل السائق ضمن الحدّ؟
\end{exercise}

\begin{exercise}[$\star\star\star$]\label{exo:g12:kinematics-2d:14}
يتسارع مترو من السكون بمقدار \qty{1.2}{m/s^2} حتى \qty{20}{m/s}،
ثم ينساب \qty{40}{s}، ثم يكبح بمقدار \qty{1.5}{m/s^2} إلى التوقف.
احسب مدة كل طور وطوله، والمسافة الكلية و
\omterm{def:g10:relative-motion:speed}{السرعة المتوسطة}. وارسم $v(t)$ وتحقّق من المسافة مساحةً.
\end{exercise}

\begin{exercise}[$\star\star\star$]\label{exo:g12:kinematics-2d:15}
\omterm{def:g10:universal-gravitation:satellite}{تابع} ثابت بالنسبة إلى الأرض يدور على نصف قطر \qty{4.22e7}{m} من
مركز الأرض \omterm{def:g10:signals-and-waves:period}{بدور} \qty{86164}{s}. احسب سرعته وتسارعه
المركزي. وما الذي يوفّره، ولماذا يتعلّق \omterm{def:g10:universal-gravitation:satellite}{التابع}
فوق نقطة استوائية ثابتة واحدة (\cref{ch:g12:satellites-kepler})؟
\end{exercise}

\section{مسألة: تقرير الحادث}

\begin{problem}[{مسألة نهاية الأسبوع --- تقرير الحادث: فيلم كاميرا
لوحة قيادة يُقرأ صورةً صورة، وقانون كبح يُعاد بناؤه من آثار الانزلاق، و
دوّار يشهد على التماسك، وحكم بالكيلومترات
في الساعة}]
\label{pb:g12:kinematics-2d:1}
تصطدم سيارة بشاحنة صغيرة خرجت من شارع جانبي؛ ولا مصاب،
لكن شركات التأمين تختلف. ولدى المحقق تسجيل كاميرا علوية
(صورة كل \qty{0.50}{s})، وآثار انزلاق مقيسة، و
كاميرا لوحة قيادة السيارة. وحدّ \omterm{def:g12:kinematics-2d:velocity}{السرعة} \qty{50}{km/h}.

\medskip
\noindent\textbf{الجزء I --- قراءة الفيلم.}
تعطي الصور \omterm{def:g10:orders-of-magnitude:prefixes}{السابقة} لأي ردّ فعل، على امتداد الطريق المستقيم:
\begin{center}\small
\begin{tabular}{lccccc}
$t$ (\unit{s}) & 0.00 & 0.50 & 1.00 & 1.50 & 2.00\\
$x$ (\unit{m}) & 0.0 & 9.0 & 18.0 & 27.0 & 36.0
\end{tabular}
\end{center}
\begin{enumerate}
  \item لماذا يجب أن يسمّي التقرير أولًا مرجعًا ومحاور؟ وأيّها
        مستعمل هنا؟
  \item احسب \omterm{def:g10:relative-motion:speed}{السرعة المتوسطة} على كل من الفترات الأربع.
  \item وأيّ صنف من الحركات هذا؟ برّر من الجدول.
  \item أعطِ \omterm{def:g12:kinematics-2d:velocity}{سرعة} السيارة $v_0$ \omterm{def:g10:orders-of-magnitude:unit}{بوحدة} \unit{m/s} وبوحدة \unit{km/h}.
  \item وما التسارع خلال هاتين الثانيتين؟
\end{enumerate}

\medskip
\noindent\textbf{الجزء II --- آثار الانزلاق.}
تعطي اختبارات الإطارات على هذا الإسفلت الجافّ تباطؤ كبح ثابتًا
$a = \qty{6.0}{m/s^2}$؛ وطول آثار الانزلاق $d_b = \qty{27}{m}$.
خذ $t = 0$، و$x = 0$ عند بدء الكبح، \omterm{def:g12:kinematics-2d:velocity}{والسرعة} المجهولة $v_B$.
\begin{enumerate}[resume]
  \item بالدوال الأصلية، اكتب $v(t)$ و$x(t)$ أثناء الكبح.
  \item عبّر عن زمن التوقف $t_s$؛ وبيّن أن $d_b = v_B^2/(2a)$.
  \item استنتج $v_B$ من آثار الانزلاق. أمتسق مع الجزء I؟
  \item وزمن ردّ الفعل \qty{1.0}{s}: احسب مسافة
        ردّ الفعل ومسافة التوقف الكلية.
  \item اشرح، انطلاقًا من السؤال 7، لماذا تربّع مضاعفةُ \omterm{def:g12:kinematics-2d:velocity}{السرعة}
        مسافةَ الكبح؛ واحسب $d_b$ عند \qty{50}{km/h} بالضبط.
\end{enumerate}

\medskip
\noindent\textbf{الجزء III --- الدوّار يشهد.}
قبل ذلك، تُظهر كاميرا لوحة القيادة السيارة تدور في مسرب دوّار
نصف قطره $R = \qty{12.5}{m}$ \omterm{def:g12:kinematics-2d:velocity}{بسرعة} ثابتة، وربع دورة في
\qty{2.5}{s}. ويقدر الطريق أن يوفّر $a_n = \mu g$ على الأكثر، بأخذ
$\mu = 0.70$.
\begin{enumerate}[resume]
  \item \omterm{def:g12:kinematics-2d:velocity}{السرعة} ثابتة هناك --- فهل \omterm{def:g12:kinematics-2d:velocity}{متجهة السرعة} ثابتة؟ وهل السيارة
        متسارعة؟ اشرح.
  \item احسب \omterm{def:g12:kinematics-2d:velocity}{سرعة} السيارة في الدوّار.
  \item احسب تسارعها المركزي (ناظمًا واتجاهًا).
  \item احسب أقصى \omterm{def:g12:kinematics-2d:velocity}{سرعة} بلا انزلاق. فهل كانت السيارة ضمنها؟
  \item أعطِ \omterm{def:g10:signals-and-waves:period}{دور} دورة كاملة وتواترها عند تلك \omterm{def:g12:kinematics-2d:velocity}{السرعة}.
\end{enumerate}

\medskip
\noindent\textbf{الجزء IV --- الحكم.}
ظهرت الشاحنة على بعد $D = \qty{41}{m}$ أمام السيارة في اللحظة
التي بدأ فيها السائق ردّ فعله.
\begin{enumerate}[resume]
  \item بكم، \omterm{def:g10:orders-of-magnitude:unit}{بوحدة} \unit{km/h} وبالنسبة المئوية، تجاوزت \omterm{def:g12:kinematics-2d:velocity}{السرعة}
        قبل الكبح الحدَّ؟
  \item احسب \omterm{def:g12:kinematics-2d:velocity}{سرعة} السيارة عند الاصطدام.
  \item أعِد حساب التوقف عند \qty{50}{km/h} بالضبط (بردّ الفعل نفسه،
        والكبح نفسه): فهل تتوقف السيارة، وبأي هامش؟
  \item أعطِ الصيغة العامة لمسافة التوقف $d(v)$؛ ولماذا تدخل
        \omterm{def:g12:kinematics-2d:velocity}{السرعة} مرتين --- مرةً خطيًّا ومرةً مربَّعة؟
  \item الحكم، في جملة واحدة بأرقامها: أيّ \omterm{def:g12:kinematics-2d:velocity}{سرعة} كان
        السائق يقود بها، وهل كان الحادث سيقع عند الحدّ؟
\end{enumerate}
\end{problem}
